
\SGARefTarget{sga7.sga7i.evii.n3-4-8}{3.4.8}. Consider now an exact sequence of chain complexes
\begin{equation}\tag{\SGARefTarget{sga7.sga7i.evii.n3-4-8-1}{3.4.8.1}}
 0\longrightarrow L_\bullet\xrightarrow{f}L'_\bullet
 \xrightarrow{g}L''_\bullet\longrightarrow0.
\end{equation}
It gives functors of the form (\SGARefLink{sga7.sga7i.evii.n3-4-3}{3.4.3}),
\[
 \underline\Psi_{L''_\bullet}\xrightarrow{g^*}
 \underline\Psi_{L'_\bullet}\xrightarrow{f^*}
 \underline\Psi_{L_\bullet}.
\]
The transitivity property implicit in functors of the form (\SGARefLink{sga7.sga7i.evii.n3-4-3}{3.4.3})
gives an isomorphism of functors
\[
 f^*g^*\simeq(gf)^*=(0_{L_\bullet\to L''_\bullet})^*.
\]
The last member is evidently canonically isomorphic to the
cocartesian zero functor from $\underline\Psi_{L''_\bullet}$ to
$\underline\Psi_{L_\bullet}$.  The preceding isomorphism therefore
defines a canonical functor
\begin{equation}\tag{\SGARefTarget{sga7.sga7i.evii.n3-4-8-2}{3.4.8.2}}
 \underline\Psi_{L''_\bullet}\longrightarrow
 \underline\Psi_{[L_\bullet\longrightarrow L'_\bullet]}.
\end{equation}
Here $\underline\Psi_{[L_\bullet\longrightarrow L'_\bullet]}$
denotes the category of pairs $(X',t)$ in which $X'$ is an object of
$\underline\Psi_{L'_\bullet}$ and $t$ is a trivialization of its
inverse image $f^*(X')$, that is, an isomorphism
\[
 f^*(X')\simeq0_{L_\bullet,A},
\]
where $0_{L_\bullet,A}$ denotes the zero object of
$\underline\Psi_{L_\bullet}(A)$ and $A$ is the object of
$\underline C$ underlying $X'$.  We claim that \underline{the
functor} (\SGARefLink{sga7.sga7i.evii.n3-4-8-2}{3.4.8.2}) \underline{is an equivalence of categories}.
Together with the left exactness in $A$ established in \SGARefLink{sga7.sga7i.evii.n3-3-1}{3.3.1}, this
completes the exactness property in $L_\bullet$ of the expression
$\underline\Psi_{L_\bullet}(A)$.

Let us sketch the proof.  To give an object of the target of
(\SGARefLink{sga7.sga7i.evii.n3-4-8-2}{3.4.8.2}) is to give an extension
\[
 0\longrightarrow A\xrightarrow{i'}X'\xrightarrow{j'}L'_0
 \longrightarrow0
\]
of $L'_0$ by an object $A$ of $\underline C$, together with a
morphism $\alpha':L'_1\longrightarrow X'$ satisfying
$j'\alpha'=d'_0$ and $\alpha'd'_1=0$, and a morphism
$t:L_0\longrightarrow X'$ satisfying $j't=f_0$ and
$td_0=\alpha'f_1$.  These data are displayed in the following
diagram:
\begin{center}
\begin{tikzpicture}[
  x=1cm,y=1cm,>=Latex,
  every node/.style={font=\small,inner sep=1.2pt}
]
  \node (z00) at (0,0) {$0$};
  \node (l0)  at (1.3,0) {$L_0$};
  \node (lp0) at (2.8,0) {$L'_0$};
  \node (lpp0) at (4.3,0) {$L''_0$};
  \node (z05) at (5.5,0) {$0$};

  \node (z10) at (0,-1.35) {$0$};
  \node (l1)  at (1.3,-1.35) {$L_1$};
  \node (lp1) at (2.8,-1.35) {$L'_1$};
  \node (lpp1) at (4.3,-1.35) {$L''_1$};
  \node (z15) at (5.5,-1.35) {$0$};

  \node (z20) at (0,-2.70) {$0$};
  \node (l2)  at (1.3,-2.70) {$L_2$};
  \node (lp2) at (2.8,-2.70) {$L'_2$};
  \node (lpp2) at (4.3,-2.70) {$L''_2$};
  \node (z25) at (5.5,-2.70) {$0$};

  \node (xp) at (3.25,1.25) {$X'$};
  \node (ap) at (4.0,2.25) {$A$};
  \node (zap) at (4.75,3.25) {$0$};
  \node (xpp) at (5.85,1.25) {$X''$};
  \node (app) at (6.60,2.25) {$A$};
  \node (zapp) at (7.35,3.25) {$0$};

  \draw[->] (z00)--(l0);
  \draw[->] (l0)--node[above]{$f_0$}(lp0);
  \draw[->] (lp0)--node[above]{$g_0$}(lpp0);
  \draw[->] (lpp0)--(z05);
  \draw[->] (z10)--(l1);
  \draw[->] (l1)--node[above]{$f_1$}(lp1);
  \draw[->] (lp1)--node[above]{$g_1$}(lpp1);
  \draw[->] (lpp1)--(z15);
  \draw[->] (z20)--(l2);
  \draw[->] (l2)--(lp2);
  \draw[->] (lp2)--(lpp2);
  \draw[->] (lpp2)--(z25);

  \draw[->] (l1)--node[right]{$d_0$}(l0);
  \draw[->] (lp1)--node[right]{$d'_0$}(lp0);
  \draw[->] (lpp1)--node[right]{$d''_0$}(lpp0);
  \draw[->] (l2)--(l1);
  \draw[->] (lp2)--node[right]{$d'_1$}(lp1);
  \draw[->] (lpp2)--node[right]{$d''_1$}(lpp1);

  \draw[->] (zap)--(ap);
  \draw[->] (ap)--node[left]{$i'$}(xp);
  \draw[->,dashed] (zapp)--(app);
  \draw[->,dashed] (app)--node[left]{$i''$}(xpp);
  \draw[->] (l0) to[bend left=22] node[above]{$t$} (xp);
  \draw[->] (xp)--node[left]{$j'$}(lp0);
  \draw[->] (lp1)--node[pos=.82,right]{$\alpha'$}(xp);
  \draw[->,dashed] (xpp)--node[above]{$j''$}(lpp0);
  \draw[->,dashed] (lpp1) to[bend right=25]
    node[pos=.82,right]{$\alpha''$} (xpp);

  \node[anchor=west] at (7.55,-0.70) {$
    \left\{
    \begin{aligned}
      j'\alpha'&=d'_0,\\
      \alpha'd'_1&=0,\\
      j't&=f_0,\\
      td_0&=\alpha'f_1.
    \end{aligned}
    \right.$};
\end{tikzpicture}
\end{center}

The familiar left exactness, in the first argument, of extensions
of an object $L$ by an object $A$, applied to the first row of the
diagram, shows that a morphism $t$ satisfying $j't=f_0$ allows one
to interpret $X'$ as the inverse image of an extension $X''$ of
$L''_0$ by $A$.  Under this interpretation, $t$ corresponds to the
natural trivialization of the inverse image of that extension along
$g_0f_0=0$.

The morphism $\alpha'$ satisfying $j'\alpha'=d'_0$ and
$\alpha'f_1=td_0$ is interpreted by essentially the same dictionary,
now applied to the second row of the diagram, as arising from a
morphism $\alpha'':L''_1\longrightarrow X''$ satisfying
$j''\alpha''=d''_0$.  Here $\alpha'$ and $\alpha''$ are regarded as
trivializations of the inverse images of $X'$ and $X''$ along
$d'_0$ and $d''_0$, respectively.  Finally, the relation
$\alpha'd'_1=0$ is then equivalent to $\alpha''d''_1=0$.  Thus the
data $(X',\alpha',t)$, an object of
$\underline\Psi_{[L_\bullet\longrightarrow L'_\bullet]}$, are
equivalent to the data $(X'',\alpha'')$, an object of
$\underline\Psi_{L''_\bullet}$.  This proves, in substance, that
(\SGARefLink{sga7.sga7i.evii.n3-4-8-2}{3.4.8.2}) is indeed an equivalence of categories.

\SGARefTarget{sga7.sga7i.evii.n3-4-9}{3.4.9}. Continue to consider the exact sequence (\SGARefLink{sga7.sga7i.evii.n3-4-8-1}{3.4.8.1}), and the
exact sequence of $\Ext^i$ groups that it defines:
\begin{multline}\tag{\SGARefTarget{sga7.sga7i.evii.n3-4-9-1}{3.4.9.1}}
 0\longrightarrow\Ext^0(L''_\bullet,A)
 \xrightarrow{g^0}\Ext^0(L'_\bullet,A)
 \xrightarrow{f^0}\Ext^0(L_\bullet,A)
 \xrightarrow{\partial^0}\Ext^1(L''_\bullet,A)\\
 \xrightarrow{g^1}\Ext^1(L'_\bullet,A)
 \xrightarrow{f^1}\Ext^1(L_\bullet,A)\longrightarrow\cdots.
\end{multline}
The canonical isomorphisms of the forms (\SGARefLink{sga7.sga7i.evii.n3-1-8}{3.1.8}) and (\SGARefLink{sga7.sga7i.evii.n3-2-4}{3.2.4}) give a
geometric interpretation of the first six terms of this exact
sequence in terms of the categories $\underline\Psi_{L_\bullet}$,
$\underline\Psi_{L'_\bullet}$, and
$\underline\Psi_{L''_\bullet}$.  Compatibility (\SGARefLink{sga7.sga7i.evii.n3-4-6}{3.4.6}) likewise
gives a geometric interpretation of the arrows
$f^0,g^0,f^1,g^1$.  We shall now give such an interpretation of
$\partial^0$, using the result just proved that (\SGARefLink{sga7.sga7i.evii.n3-4-8-2}{3.4.8.2}) is an
equivalence of categories.

Let
\[
 u\in\Ext^0(L_\bullet,A),
 \qquad\text{i.e.}\qquad
 u:L_0\longrightarrow A,\quad ud_0=0.
\]
Consider the object $(0,\overline u)$ of
$\underline\Psi_{[L_\bullet\to L'_\bullet]}$.  Here $0$ is the
zero object of $\underline\Psi_{L'_\bullet}$, whose inverse image
along $f_0$ is therefore identified with the zero object of
$\underline\Psi_{L_\bullet}$, and $\overline u$ is the automorphism
of the latter defined by $u$ through the canonical isomorphism
(\SGARefLink{sga7.sga7i.evii.n3-1-8}{3.1.8}).  Since (\SGARefLink{sga7.sga7i.evii.n3-4-8-2}{3.4.8.2}) is an equivalence, $(0,\overline u)$ comes
from an object of $\underline\Psi_{L''_\bullet}(A)$, determined up
to unique isomorphism.  Its isomorphism class is an element of
\[
 \Psi^1_{L''_\bullet}(A)\simeq\Ext^1(L''_\bullet,A).
\]
\underline{We claim that this element is precisely}
$\partial^0(u)$.

We prove this compatibility by reducing it to the analogous one in
[4, Cor.~5.5], where the sign must be corrected as indicated in
3.3.6.3.  The argument also gives another proof of the equivalence
(\SGARefLink{sga7.sga7i.evii.n3-4-8-2}{3.4.8.2}), by reducing it to [4, Prop.~\SGARefLink{sga7.sga7i.evii.n2-7}{2.7}].  Consider the opposite
category $\underline C^{\mathrm o}$ and interpret the construction
of $\underline\Psi_{L_\bullet}(A)$, for fixed $A$, together with
the constructions of \SGARefLink{sga7.sga7i.evii.n3-4-8}{3.4.8} and the present section, in terms of
$\underline C^{\mathrm o}$ and the category $\underline E$ cofibred
over $\underline C^{\mathrm o}$ whose objects are extensions of the
fixed object $A$ of $\underline C^{\mathrm o}$ by a variable object
$L$ of $\underline C^{\mathrm o}$.  This is an additive cofibred
category over $\underline C^{\mathrm o}$, and is even left exact
there (in the sense of [4, \SGARefLink{sga7.sga7i.evii.n2-2}{2.2}]).

A chain complex $L_\bullet$ of $\underline C$ may be interpreted as
a cochain complex of $\underline C^{\mathrm o}$.  Under this
interpretation, the construction of $\underline\Psi_{L_\bullet}(A)$
coincides with the opposite of the category
$\underline{\widehat E}(L_\bullet)$ in the sense of [4, 1.8].  The
functor (\SGARefLink{sga7.sga7i.evii.n3-4-8-2}{3.4.8.2}) is then identified with the opposite of the functor
of the form [4, (2.1.2)], with $E$ replaced by its extension
$\underline{\widehat E}$ to the category of cochain complexes.
Consequently, the assertion that this functor is an equivalence is
the special case [4, \SGARefLink{sga7.sga7i.evii.n2-7}{2.7}, a)$\Rightarrow$d)].

On the other hand, the exact sequence (\SGARefLink{sga7.sga7i.evii.n3-4-9-1}{3.4.9.1}) is identified with
the exact sequence of the groups $\Ext^i(A,L''_\bullet)$,
$\Ext^i(A,L'_\bullet)$, and $\Ext^i(A,L_\bullet)$, apart from a
factor $(-1)^i$ on $\partial^i$ (Verdier dixit).  Similarly, the
isomorphisms (\SGARefLink{sga7.sga7i.evii.n3-1-8}{3.1.8}) and (\SGARefLink{sga7.sga7i.evii.n3-2-4}{3.2.4}) are identified with the analogous
isomorphisms (3.11.5) and (3.11.4).  The asserted compatibility is
therefore the special case of [4, Cor.~5.5] corrected as above.

\underline{Remark} \SGARefTarget{sga7.sga7i.evii.n3-4-9-2}{3.4.9.2}. The editor, being very afraid of signs,
confesses that he does not have complete confidence in the preceding
proof: he fears that signs may have been omitted in passing from
$C$ to $C^{\mathrm o}$.  He recommends that a courageous reader
verify the asserted compatibility directly (or, respectively, its
opposite).  Fortunately, these sign questions will play no role in
the following expos\'es.
